% Draft unit for YM-000026.  Parent supplies theorem environments/macros.
\section{The annular Ward field for the selected stress tensor}
\label{sec:ex-lattice-actual-stress-annulus}

This section proves the shrinking-annulus implication used by the public
stress-tensor route.  Given the selected shell-node receipt, a finite linear
combination of scale-covariant insertions that annihilates every field at one
grade has vanishing flux on shrinking windows.  Finite-dimensional linear
algebra then represents those fluxes by one field at that grade.  The
production of the shell-node receipt is a separate quantitative obligation;
it is not proved by the conditional argument below.

Write $a_i$ for the physical spacing at level $i$ and
$\lambda_i=\log(a_i^{-1})\to\infty$ for its logarithmic scale; thus
$a_i=e^{-\lambda_i}$.  Write $B_i$ for its support radius, and
\[
  r_m=128 B_{m+1}^{2}a_{m+1}.
\]
The scale note used below says $r_m\to0$.  Let $\mathcal F_d$ denote the
complex field space at grade $d$, and let $\tau_{I,D}:\mathcal F_D\to\mathbb C$
be the tester associated with a watch $I$.  We write
\[
  \mathcal A_r=\{y\in\mathbb R^4:r/4<|y|<2r\},\qquad
  \mathcal C_r=\{y\in\mathbb R^4:r/16<|y|<4r\}.
\]
The wider collar $\mathcal C_r$ is where the finite-volume shell correction is
constructed; the desired annihilation is on $\mathcal A_r$.

Concretely, $\tau_{I,D}(z)$ is the insertion functional obtained by placing
the grade-$D$ field $z$ at the distinguished point of $I$ and retaining its
other points and color labels as spectators.  The flux functional is the
window integral of the stress insertion against the window gradient and the
Killing field:
\[
 \operatorname{Flux}_{J}(\chi,U,I,v)
 =-\int_{\mathbb R^4}\sum_{\mu,\nu=1}^4
   \partial_\mu\chi(x)\,U_\nu(x-x_I)\,
   (I)\!\operatorname{asked}(J_{\mu\nu},v)(x)\,dx,
\]
where $x_I$ is the distinguished point of the watch; all fluxes are thus
written in coordinates centered at that point.  The sign is fixed by the
flux convention in the stress conservation law.
All insertions and integrals here are the selected-table quantities; the
shell argument uses their stated continuity and local conservation
properties.  The approximating counterterms below are the finite-volume
objects.

For a finite-volume cell system $\mathcal T$ and coupling $\beta$, the
root-test of a configuration observable $F$ is
\[
 \operatorname{RT}_{\beta,\mathcal T}(F)
 =\left|\mathbb E_\beta\prod_{u\in\mathcal T}F_u\right|^{1/|\mathcal T|},
\]
where $F_u$ is the prescribed translate of $F$ at cell $u$, with the source
reflection conjugation when that cell is reflected.  The label norm
$\|z\|_D$ below is the fixed coefficient norm of the grade-$D$ label space.

For the exact shell costs, put
\[
 c_m=e^{\lambda_m/32},\quad
 S(p_0,k)=\exp\!\left((a_{p_0}/a_k)^4\,\mathsf T_k+5\right),\quad
 X_m=\exp\!\sqrt{\widehat\beta_{m+1}},
\]
where $\mathsf T_k$ is the treasury quantity at level $k$.  Let $V_m>0$ be
the product of the two selected channel payments at levels $(j_m,m+1)$, and
let $n_m$ be the number of columns at level $m$.  The counterterm budget is
\[
 e_m=32(1+n_m)a_{m+1}^{D}X_m/V_m.
\]
The quota conversion supplies a finite factor $Q_0$ for the fixed
grade/threshold, so the correction label has coefficient norm at most
$2Q_0e_m$ at annulus level $m$.  The root-test cost, in contrast, is the
decaying quantity $2S(p_0,k)c_m^{-1}$.

\subsection{The shell-node hypothesis}

Fix grades $d,D,h\in\mathbb N$ with
\begin{equation}\label{eq:ex-lattice-stress-grade}
  D+h=d+4.
\end{equation}
We name the exact shell-node input consumed below.  This definition does not
assert that the receipt has been produced from the public history.

\begin{definition}[Shell-node receipt]
\label{def:ex-lattice-shell-receipt}
For fixed grades $d,D,h$ satisfying \eqref{eq:ex-lattice-stress-grade}, a
fixed family and selected history have the shell-node receipt if there are
$t>1$ and levels $j_m$, chosen before the opened levels, labels, or extraction
order, such that
\[
 j_m\longrightarrow\infty,\qquad
 \frac{\lambda_{j_m}}{\lambda_m}\longrightarrow t,
\]
and the scale note has the consequence
\[
                         \frac{\log r_m}{\lambda_m}\longrightarrow-1.
\]
For every fixed opened reference level $p_0$ and level $k$ with $p_0$ open at $k$ and
with $k$ at least the reference threshold for grade $D$, every $v\in\mathcal F_d$,
$T\in\mathcal F_4$, and extraction order $\mathfrak f$, there is $m_0$ such
that for every $m\geq m_0$ and every fixed integer lattice anchor
$x\in\mathrm{Cover}=\mathbb Z^4$ whose continuum location at level $j_m$
lies in $\mathcal C_{r_m}$, there are a label sequence
$z_x: \mathbb N\times\mathbb N\to\mathcal F_D$ and a constant $A'_m$ for which
$\mathfrak f$-almost every $(N,q)$ has $p_0\leq N$ and
\[
  \|z_x(N,q)\|_D\leq2Q_0e_m,\qquad
  \operatorname{RT}_{\beta,\mathcal T}(\text{fragment}(z_x(N,q)))
       \leq A'_m:=2S(p_0,k)c_m^{-1}.
\]
The late filter set and sequence may depend on $x$; the numeric coefficient
and root-test bounds do not.  The factor $Q_0$ is independent of volume,
period, and collar point; the coefficient bound may depend on $m$ through
$e_m$.  The finite extraction order may encode any fixed finite family of
watches; the receipt is available for every such family, with its tester costs
fixed before the eventual volume/period choice.
The associated finite-volume fragment is the product of the grade-four
insertion at the collar point and the grade-$d$ insertion at the selected
point, multiplied by $r_m^h$, minus the grade-$D$ insertion $z_x$ at the
selected point.  It lies in the prescribed chamber and has root-test norm
bounded by the displayed cost.  No common late filter over the entire (growing)
set of collar grid points is asserted.
\end{definition}

Here ``eventually'' is on the extraction filter over the volume and period
variables, not a claim at each finite volume.  The source construction has two
stages.  First, a two-channel counterterm construction builds the fragment.
At the collar site $x$, its two factors are the
base insertion and the rescaled grade-four insertion; their product is
corrected by one grade-$D$ label.  The grade identity \eqref{eq:ex-lattice-stress-grade}
is exactly what makes that correction have the requested degree.  The paid
two-component insertion bounds its root-test cost, while the independent
quota estimate bounds the correction label.  Second, the corner-to-shell
argument supplies common numeric bounds for all collar anchors.  The late
extraction event remains pointwise in the anchor; for a fixed interpolation
stencil, one obtains a common event by intersecting the finitely many
pointwise late events.  This does not assert one event for the entire growing
collar.

The source route proposes the following numerical gates for this input.  The
economics choice is
\[
  t=1+\frac{1}{64(d+12)},
\]
with a subsequence $j_m$ for which $\lambda_{j_m}/\lambda_m\to t$.  The price
estimate supplies, eventually, positivity and a small two-channel payment
budget, together with the grade-$D$ correction-price bound.  In the source
normalization its exponential denominator is $c_m=\exp(\lambda_m/32)$; at fixed
opened $p_0,k$, $S(p_0,k)$ is fixed, so the root cost tends to zero.  The separate
lightness estimate supplies both
\[
  32r_m\leq a_{p_0},\qquad
  a_{m+1}\bigl(B_{m+1}^2+1+2\cdot1024 B_{m+1}^2\bigr)\leq a_{p_0}
\]
eventually.  These are used, respectively, for the collar/spacing margin and
for the number of required coordinates in the opened chamber.  The corner
lemma combines them with the same-history run, throw, root-room, and
subsidized-payment bounds.  This draft does not prove that those source
receipts imply the shell-node input above; that quantitative producer remains
open.

The implication below takes the defined shell-node receipt as an input
on one fixed family.  It does not infer that receipt from the public
extractor's record or change family, history, or scale sequence between the
receipt and the stress law.

\subsection{Vanishing against a finite tester annihilator}

\begin{lemma}[shrinking-annulus annihilation]
\label{lem:ex-lattice-annulus-annihilation}
Let $S$ be the selected hierarchy, $V$ a derivative plan for $S$, and $J$ a
stress conserved with respect to $V$ and the spin action of $S$.  Assume the
lattice configurations used in the shell argument are dense in the separated
configuration space and that the flux and tester functionals are continuous
there.  Fix an opened reference level $p_0$ and level $k$ as in the shell-node
receipt.  Take $v\in\mathcal F_d$, and let
$U:\mathbb R^4\to\mathbb R^4$ be a smooth homogeneous Killing field of degree
$p\in\mathbb N$, so $U(ry)=r^pU(y)$ for $r>0$.  Suppose both
\begin{equation}\label{eq:ex-lattice-flux-grade}
  D+(3+p)=d+4
\end{equation}
and \eqref{eq:ex-lattice-stress-grade} hold, and the shell-node receipt
holds at $(p_0,k)$ for the finite family of watches
$I_1,\ldots,I_s$.  Fix its witnessing scale ratio $t>1$.  For each watch
$I_a$, the moving insertion is a component of the grade-four stress, so its
initial tag degree is $q_a=4$; let $c_a$ be the total cloud degree of all
spectators, including the fixed grade-$d$ insertion $v$.  Choose an integer $q_*\geq1$ so large
that
\[
  q_a+q_*+2+c_a<q_*t
  \quad (1\leq a\leq s).
\]
For $0\leq n\leq q_*$, let $H_{a,n,m}(y)$ denote the scalar insertion with
the moving stress component of $J$ at $y$ raised $n$ times, with the watch's
fixed grade-$d$ insertion $v$ and all other spectator labels held fixed.  Assume the quantitative
separated-point bounds: for every sufficiently large $m$ and every
$y\in\mathcal C_{r_m}$,
\[
 |H_{a,n,m}(y)|\leq C_a|y|^{-(q_a+n+2+c_a)},
 \qquad 0\leq n\leq q_* ,
\]
where $C_a$ is independent of $m$ and $y$. In physical coordinates,
\[
 \partial_\ell H_{a,n,m}=H_{a,n+1,m}.
\]
Thus along the source chord
$t\mapsto x+a_{j_m}t e_\ell$ the derivative is
$a_{j_m}H_{a,n+1,m}$; the $n$th chain member is
$a_{j_m}^{n}H_{a,n,m}(x+a_{j_m}t e_\ell)$.  These explicit quantitative estimates, rather than smoothness
alone, are the local regularity input used below.  If
$g_1,\ldots,g_s\in\mathbb C$ satisfy
\begin{equation}\label{eq:ex-lattice-annulus-annihilator}
  \sum_{a=1}^s g_a\tau_{I_a,D}(z)=0
  \quad\hbox{for every }z\in\mathcal F_D,
\end{equation}
then, for any smooth windows $\chi_a$ valid at the respective watch
configurations,
\begin{equation}\label{eq:ex-lattice-annulus-flux-zero}
  \sum_{a=1}^s g_a\,\operatorname{Flux}_{J}
       (\chi_a,U,I_a,v)=0.
\end{equation}
\end{lemma}

\begin{proof}
Choose a smooth radial bump $\varphi$ on $\mathbb R^4$ which equals one on
the ball of radius $1/2$ and vanishes outside the unit ball.  The selected
watch configurations have their distinguished point at the origin and all
other points separated from it.  Since $r_m\to0$, for all sufficiently large
$m$ the rescaled bump $\varphi_m(y)=\varphi(y/r_m)$ is a valid window at every
watch in the finite family, and no other insertion lies in its support.
Conservation and the window-change identity therefore let us compute each
flux using this same shrinking window.

For each pair of tensor indices $i,j$, define
\[
  W_{ij}(y)=\partial_i\varphi(y)\,U_j(y),\qquad
  G_{ij,m}(y)=W_{ij}(y)\,r_m^{3+p}
       \sum_{a=1}^s g_a\,(I_a)\!\operatorname{asked}(J_{ij},v)(r_my).
\]
Each $W_{ij}$ is integrable and supported in the unit ball.  Its derivative
vanishes near the origin, and the bump-support calculation places every point
where $W_{ij}$ is nonzero in the annulus with radii between $r_m/4$ and
$2r_m$ after scaling.  The annihilation argument below uses the shell
estimate on the slightly wider collar, so endpoint conventions do not matter.

For each $m$, the shell-node receipt has root-test tolerance
\[
  \epsilon_m=2S(p_0,k)c_m^{-1}\longrightarrow0.
\]
For this fixed $m$ the interpolation stencil uses only finitely many grid
points.  Intersect the late filter sets supplied separately for those points;
their finite intersection is still late.  The corresponding finite tuple of
correction labels lies in a bounded subset of a finite-dimensional product
space, so pass to one convergent subnet and take its coordinate limits.  This
single extraction handles every stencil point and every watch in the finite
family; no subsequence is chosen watch by watch.  The
watch-cost limits identify each limiting tester with the corresponding
rescaled product insertion, up to its paid error $A_a\epsilon_m$.  Summing
with coefficients $g_a$, the limiting grade-$D$ terms cancel by
\eqref{eq:ex-lattice-annulus-annihilator}; the remaining norm is at most
$\epsilon_m\sum_a|g_a|A_a$.  Thus the grid error tends to zero.  Only this
fixed finite family is used; no uniform bound over all watches is asserted.
The one-coordinate Taylor argument described below transfers this grid error
to the annulus.  The two grade relations give $h=3+p$.  Thus the expression
multiplied by $W_{ij}$ is precisely the scale-corrected insertion that occurs
after changing variables $x=r_my$ in the flux integral: the window gradient
contributes $r_m^{-1}$, the homogeneous field contributes $r_m^p$, and
Lebesgue measure contributes $r_m^4$.  Their product is $r_m^{3+p}$, and
$h=d+4-D=3+p$.  Hence this is exactly the factor
$r_m^h$ in the finite-tester estimate; there is no additional scaling factor.

The absolute value of each rescaled component integral is bounded by the
uniform shell error times $\int |W_{ij}|$.  There are only sixteen components,
so summing over them still tends to zero.  The change-of-variables identity
for the flux gives
\[
 \sum_a g_a\operatorname{Flux}_{J}(\varphi_m,U,I_a,v)
   =-\sum_{i,j}\int_{\mathbb R^4}G_{ij,m}(y)\,dy.
\]
The left side is independent of $m$ by conservation and the support
separation.  The right side tends to zero by the preceding integrable bound.
Therefore that fixed left side is zero.  Replacing $\varphi_m$ by $\chi_a$
uses the same window-change identity separately at each watch, proving
\eqref{eq:ex-lattice-annulus-flux-zero}.\end{proof}

The grid-to-annulus interpolation has an elementary quantitative form.  For
$q\geq0$ and $P>0$, set
\[
 A_0(P)=1,\qquad A_{q+1}(P)=1+4P A_q(P).
\]
Suppose $f_0,\ldots,f_q:[-P,P]\to\mathbb C$ are differentiable with
$f_i'=f_{i+1}$ for $i<q$, $|f_q|\leq\epsilon$ throughout the interval,
and $|f_0(x_j)|\leq\epsilon$ at $2^q$ ordered points whose consecutive
gaps are at least one.  Then
\begin{equation}\label{eq:ex-lattice-annulus-interpolation}
 \sup_{[-P,P]}|f_0|\leq2A_q(P)\epsilon.
\end{equation}
For the proof first take all functions real-valued.  Pair consecutive
nodes.  The real mean-value theorem supplies an intermediate point at which
$f_1$ is the difference quotient; its absolute value is at most
$2\epsilon$ because the gap is at least one.  The $2^{q-1}$ intermediate
points remain ordered with gaps at least one.  Induction applied to
$f_1,\ldots,f_q$ bounds $|f_1|$ by $2A_{q-1}(P)\epsilon$.  Integrating from
one original node to an arbitrary point costs at most $2P$ times that bound,
and the value at the node costs at most $\epsilon$, giving
$A_q(P)\epsilon$.  For complex functions apply this real estimate
separately to real and imaginary parts.  Their two bounds imply the stated
complex estimate, since $|z|\leq|\Re z|+|\Im z|$.

The common-limit argument applies separately at each of the finitely many
grid nodes in the interpolation stencil.  Its estimate is uniform in the
node because the shell-node receipt and the watch-cost constants are uniform on
the collar.  At those nodes the nonzero-coordinate condition places every
watch in the generic region required by its cost-limit estimate.

To pass from those grid points to an arbitrary point of the annulus, the
derivative plan supplies the chains in
\eqref{eq:ex-lattice-annulus-interpolation}.  Along a coordinate line,
differentiating a disc insertion raises its first tag by one and multiplies
by the level spacing.  Thus the $i$th chain member is the finite sum with
first tag raised $i$ times and prefactor $a_{j_m}^i$, and its derivative is
the next member.  Since $\lambda_{j_m}/\lambda_m\to t>1$ while $r_m$ has logarithmic spatial
rate $-1$, the choice of $q_*$ above gives
\[
 q_a+q_*+2+c_a<q_*t,
 \qquad\text{equivalently}\qquad
 q_*(t-1)>q_a+2+c_a.
\]
The iterated derivative estimate gives a uniform top-derivative bound of the form
\[
 C a_{j_m}^{q_*}r_m^{-(q_a+q_*+2+c_a)},
\]
where $C$ is independent of $m$.  Its logarithm divided by $\lambda_m$
tends to $-q_*t+(q_a+q_*+2+c_a)<0$, so the bound tends to zero.
For the finite-node interpolation, induct over the set $K$ of axes not yet
discretized.  On the next axis choose the nonzero shifted lattice coordinate
provided by the one-dimensional stencil; leave the other coordinates fixed.
At an intermediate stage only the already chosen coordinates are known to be
nonzero.  At the final stage all four are nonzero, so the grid-point representation
and genericity estimates apply.  Each coordinate move is at most $r_m/128$;
therefore the total displacement is at most $|K|r_m/128\le r_m/32$, within
the collar margin.  The collar-margin condition
$\operatorname{arm}(q_*)a_{j_m}\leq r_m/128$ keeps each segment in the wider
collar.  Applying \eqref{eq:ex-lattice-annulus-interpolation} successively
in the four coordinates gives a fixed finite factor, independent of $m$.
The grid error and the top derivative both tend to zero, proving uniform
convergence on the annulus.  The constants depend on the finite family, the
grade, and the chosen chain length, but not on $m$.

\subsection{One field represents every flux}

\begin{proposition}[annular flux representation]
\label{prop:ex-lattice-annular-field}
Assume the hypotheses of Lemma~\ref{lem:ex-lattice-annulus-annihilation},
with the shell-node receipt available for every finite family of watches.  For each
$v\in\mathcal F_d$ and each smooth homogeneous Killing field $U$ of degree
$p$, there exists one field $b\in\mathcal F_D$ such that for every
watch $I$, every compatible smooth window $\chi$, and every separated point
configuration,
\[
  \operatorname{Flux}_{J}(\chi,U,I,v)=\tau_{I,D}(b).
\]
The same $b$ works for all configurations, windows, and watches.
\end{proposition}

\begin{proof}
For a finite family $I_1,\ldots,I_s$, define linear maps
\[
  T:\mathcal F_D\to\mathbb C^s,
  \quad T(z)=(\tau_{I_1,D}(z),\ldots,\tau_{I_s,D}(z)),
  \qquad a=(\operatorname{Flux}_{J}(\chi_1,U,I_1,v),\ldots).
\]
Lemma~\ref{lem:ex-lattice-annulus-annihilation} says that every linear
functional annihilating $\operatorname{range}T$ also annihilates $a$.
In finite-dimensional linear algebra, a subspace is the common kernel of all
linear functionals vanishing on it.  Hence $a\in\operatorname{range}T$ and
there is $b$ with $T(b)=a$.  This is the finite-tester step; it does not presume an
infinite-dimensional Riesz representation.

Choose a finite faithful family of watches whose testers distinguish
elements of $\mathcal F_D$.  Such a family exists because the complete family
of watch evaluations separates fields, and a finite-dimensional field space
has a finite subfamily of any separating family that still separates.  Apply
the finite result to that family to obtain
$b$.  For any further watch $I$, apply the result to the union of $I$ and the
faithful family.  The shell-node receipt, assumed for every finite family of watches,
applies to this union.  It gives $b'$.  On the faithful family, $b'$ and $b$ have
the same tester values, so faithfulness implies $b'=b$.  Therefore the flux
identity holds for $I$ as well.  This argument also handles every compatible
window, since the window-change identity identifies its flux with the chosen
window value.  Finally, the lattice configurations used above are dense
among separated configurations.  Both flux and the tester/disc functional
are continuous there, so the identity extends from those configurations to
all separated point configurations.\end{proof}

\subsection{Translation and rotation Ward maps}

Apply Proposition~\ref{prop:ex-lattice-annular-field} first to each constant
Killing field $U(y)=e_\mu$, of degree $p=0$, with output grade $D=d+1$.
For every input grade $d$ and direction $\mu$ this gives a complex-linear map
$P_{d,\mu}:\mathcal F_d\to\mathcal F_{d+1}$.  Linearity follows because
both flux and the finite faithful tester map are linear; uniqueness on the
faithful family makes the representing field linear in the input.  The
translation identity identifies the front insertion of $P_{d,\mu}v$ with
the stress flux.

Next take the rotational Killing field $U(y)=Ky$, where $K$ is skew-symmetric.
It is homogeneous of degree one, so the grade relation gives $D=d$.  The same
construction yields maps $R_{d,K}:\mathcal F_d\to\mathcal F_d$ which are
complex-linear in the field and real-linear in $K$, by linearity of the flux
and uniqueness on the faithful testers.  The angular-flux identity follows
in coordinates centered at each watch's distinguished point, as in the flux
definition.  These two families are precisely the translation
and rotation Ward maps; no additional stress tensor is selected in this
step.

\begin{remark}[public same-history application]
\label{rem:ex-lattice-stress-public-binding}
The public route keeps the relevant objects on the family selected by the
run: the extraction, hierarchy, finite-volume laws, and conserved stress all
come from that same record.  The endpoint then invokes the shell construction
on this family.  The source correspondence is exact, but this unit does not
prove that the selected history supplies the shell-node receipt defined
above; the conditional Ward conclusion cannot be applied to the public stress
until that producer is discharged.
\end{remark}

The argument proves finite-tester annihilation and the unique-field
representation conditional on the shell-node receipt and the stated
regularity/conservation data.  Producing that receipt from the selected
history, including the paid local estimates and the finite-volume
counterterm, remains open in this unit.  The source crosswalk is recorded in
the companion evidence file; no published theorem is substituted for these
project-specific quantitative estimates.
